%---------------------------------------------------------------------------%
%->> Backmatter
%---------------------------------------------------------------------------%
\chapter[致谢]{致\quad 谢}\chaptermark{致\quad 谢}% syntax: \chapter[目录]{标题}\chaptermark{页眉}
%\thispagestyle{noheaderstyle}% 如果需要移除当前页的页眉
%\pagestyle{noheaderstyle}% 如果需要移除整章的页眉

值此毕业之际，回望在江门中微子实验（JUNO）组度过的这三年硕士时光，心中满是感慨与不舍。在中微子组度过了一段充实而又令人眷恋的时光。

首先，我要向我的导师丁雪峰老师致以最崇高的敬意和最衷心的感谢。 丁老师以其深厚的代码功底、扎实的物理功底以及敏锐的物理嗅觉，为我打开了粒子物理的大门。在学术指导中，丁老师从不强加定见，而是致力于培养我独立思考的能力，鼓励我从底层逻辑出发去复现、去质疑。每当我陷入错误或逻辑困局时，丁老师总能以极大的耐心包容我的笨拙，并给出举重若轻的指导。同时也感谢丁老师在百忙之中为我撰写申博推荐信。

感谢在高能物理研究所做本科毕设期间，给予我帮助的温良剑老师和李高嵩老师，感谢你们在实验器材的使用以及实验过程中遇到棘手问题时，总能及时伸出援手，你们严谨的工作作风让我深受启发。感谢李玉峰老师和占亮老师，在反应堆中微子振荡分析方面，两位老师以极高的学术造诣为我耐心解答疑惑。每一次深入的讨论，都让我对物理图像有了更清晰的认知。同时，感谢李老师和占老师为我撰写申博推荐信，给予我迈向更高学术殿堂的信心。感谢于泽源老师，感谢于老师对我学业的关注，为我写推荐信。

感谢张涵师兄给予我极大的帮助。在反应堆振荡分析和复杂的统计学问题上，张师兄总是耐心细致地为我答疑解惑，在此深表谢意，与张涵师兄的合作让人。感谢中微子组大家庭里的每一位师兄师弟。 实验室里的头脑风暴，食堂里的学术闲谈，以及组内互帮互助的氛围，让原本枯燥的科研生活变得生动有趣。这种并肩作战的同袍之谊，是我硕士生涯中最珍贵的记忆，也是我最舍不得离开这个集体的理由。

最后，我要将最深的情感献给我的父母。 感谢你们为我营造了一个充满爱的家庭环境，让我能够拥有积极乐观的心态去面对生活中的挑战。是你们的理解与包容，让我能心无旁骛地追求科学梦想，成为一个温暖且坚定的人。

纸短情长，三年的硕士生涯即将画上句点。虽有万般不舍，但我深知，这一站的终点正是下一站的起点。


\rightline{2026年6月}
\chapter{作者简历及攻读学位期间发表的学术论文与其他相关学术成果}

\section*{作者简历：}
2019年9月——2023年6月，在中北大学半导体与物理学院获得学士学位。

2023年9月——2026年6月，在中国科学院高能物理研究所攻读硕士学位。

\section*{已发表（或正式接受）的学术论文：}

{
\setlist[enumerate]{}% restore default behavior
\begin{enumerate}[nosep]
    \item Xue Jingqin, Han Zhang, Ding Xuefeng, et al. “High-performance statistical methods for reactor neutrino oscillations,” Eur. Phys. J. C 85, 1420 (2025). DOI: https://doi.org/10.1140/epjc/s10052-025-15164-z.
\end{enumerate}
}


\section*{参加的研究项目及获奖情况：}

\begin{enumerate}[nosep]
    \item 2024年获中国科学院大学“三好学生”荣誉称号
    \item 2025年获国家奖学金
    \item 2025年获所长优秀奖学金
\end{enumerate}

\cleardoublepage[plain]% 让文档总是结束于偶数页，可根据需要设定页眉页脚样式，如 [noheaderstyle]
%---------------------------------------------------------------------------%
